\medskip

\textbf{Partie A}

\medskip

Le but de la partie A est d'étudier le comportement de la suite $\left(u_n\right)$ définie par $u_0 =  0,3$ et par la relation de récurrence, pour tout entier naturel $n$ :
\[u_{n+1} =  2u_n\left(1 - u_n\right).\]

Cette relation de récurrence s'écrit $u_{n+1} =  f\left(u_n\right)$, où $f$ est la fonction définie sur $\R$ par : 
\[f(x) = 2x\left (1 - x\right ).\]

\smallskip
\begin{enumerate}
\item Démontrer que la fonction $f$ est strictement croissante sur l'intervalle $\left[0~;~\dfrac12\right]$.
\item On admet que pour tout entier naturel $n$, $0 \leqslant u_n \leqslant \dfrac12$.

Calculer $u_1$ puis effectuer un raisonnement par récurrence pour démontrer que pour tout rentier naturel $n,\: u_n \leqslant u_{n+1}$.
\item En déduire que la suite $\left(u_n\right)$ est convergente.
\item Justifier que la limite de la suite $\left(u_n\right)$ est égale à $\dfrac12$.
\end{enumerate}

\bigskip

\textbf{Partie B}

\medskip

Le but de cette partie est d'étudier un modèle d'évolution d'une population.

En 2022, cette population compte \np{3000} individus.

\smallskip

On note $P_n$ l'effectif en milliers de la population l'année $2022 + n$. Ainsi $P_0 = 3$.

Selon un modèle inspiré du modèle de Verhulst, mathématicien belge du XIX\up{e} siècle, on considère que, pour tout entier naturel $n$ :

\begin{center}$P_{n+1} -  P_n = P_n\left(1 -  b \times P_n\right)$, où $b$ est un réel strictement positif.\end{center}

Le réel $b$ est un facteur de freinage qui permet de tenir compte du caractère limité des
ressources du milieu dans lequel évoluent ces individus.

\medskip

\begin{enumerate}
\item Dans cette question $b  = 0$.
	\begin{enumerate}
		\item Justifier que la suite $\left(P_n\right)$ est une suite géométrique dont on précisera la raison.
		\item Déterminer la limite de $P_n$.
	\end{enumerate}
\item Dans cette question $b = 0,2$.
	\begin{enumerate}
		\item Pour tout entier naturel $n$, on pose $v_n = 0,1 \times P_n$.

Calculer $v_0$ puis montrer que, pour tout entier naturel $n$,\:$v_{n+1} = 2v_n\left(1 - v_n \right)$.
		\item Dans ce modèle, justifier que la population se stabilisera autour d'une valeur que l'on précisera.
	\end{enumerate}
\end{enumerate}
